\section{Sheaf conventions and the relative Hecke system}
\label{sec:foundations}

The argument will specialize sheaves twice and will use Betti sheaves in the
intermediate characteristic-zero fiber.  We first specify the sheaf categories
and the normalization in which these operations will be compared.  In
particular, the moving points of a Hecke correspondence contribute no shift to
the eigenvalue identity.

\subsection{Local constructibility and flag levels}

Write $E=\overline{\mathbb Q}_{\ell}$.  For a smooth algebraic stack $\cB$
locally of finite type over an algebraically closed field of characteristic
different from $\ell$, let $D_{\mathrm{lcc}}(\cB,E)$ denote the category of
geometric adic complexes whose pullback to every smooth finite-type scheme
chart is bounded and constructible.  The abbreviation ``lcc'' will always
mean \emph{locally bounded constructible}, not locally constant.  Bounds may
depend on the chart.  An object $F$ is perverse if
\[
                         f^*F[d]\in\operatorname{Perv}(D,E)
\]
for every smooth chart $f:D\to\cB$ of constant relative dimension $d$.
Charts of nonconstant relative dimension are treated componentwise.  This
defines the usual perverse $t$-structure locally and hence its truncations
and perverse cohomology on $D_{\mathrm{lcc}}$.  Neither this definition nor
the word ``perverse'' imposes a support condition on the set of
Harder--Narasimhan strata.

We use ordinary pullbacks and ordinary geometric stalks.  For a closed conic
subset $C\subset T^*\cB$, its pullback to a smooth chart is
\[
 f^{\circ}C
   =\operatorname{image}\bigl(D\mathbin{\times}_{\cB}C
                                      \longrightarrow T^*D\bigr).
\]
Thus $\SS(F)\subset C$ means $\SS(f^*F)\subset f^{\circ}C$ on every such
chart.  Shifting a complex does not change this condition.  The singular
support used here is the geometric adic singular support; existence and the
function-test characterization for bounded constructible $E$-complexes are
as in \cite[\S1.5, Theorem~(vii)]{Barrett}.

Fix a maximal torus in $B$ and identify it with $T=B/N_B$, where
$N_B=R_u(B)$.  For a smooth pointed projective curve $(X,x)$ put
\[
 \cA=\Bun_{G,B,x}(X),\qquad
 \cA^+=\Bun_{G,N_B,x}(X).
\]
An enhanced flag is an $N_B$-reduction of the fiber at $x$; forgetting the
enhancement gives a representable $T$-torsor
\begin{equation}\label{eq:found-level-torsor}
                             \pi:\cA^+\longrightarrow\cA.
\end{equation}
We also allow unlevelled bundles.  For two pointed curves with identified
flag tori, and a fixed choice of identification $\tau:T_1\simeq T_2$, write
\begin{equation}\label{eq:found-node-quotient}
 \cB_{12}=\bigl[(\cA_1^+\times\cA_2^+)/T_{\Delta}\bigr],
 \qquad T_{\Delta}=\{(t,\tau(t)):t\in T_1\}.
\end{equation}
The Borels on the two curves may be opposite.  The branch convention fixes
$\tau$ once and for all.  The map from $\cB_{12}$ to
$\cA_1\times\cA_2$ is a torsor under $(T_1\times T_2)/T_{\Delta}$.
An unramified Hecke modification on either punctured component transports
the level data and acts on the corresponding factor of this presentation.
Section~\ref{sec:geometry} describes the cotangent spaces and the
generic-nilpotence cone for all these stacks.

\subsection{Geometric specialization}

Let $S=\Spec R$ be an excellent complete strictly henselian discrete
valuation trait, with algebraically closed residue field $k$, and with
$\ell$ invertible in $R$.  Fix an algebraic closure $\overline K$ of its
fraction field.  Write $\bar\eta=\Spec\overline K$ and $s=\Spec k$.
For a finite-type $S$-scheme $Y$ we use geometric nearby cycles
\[
                \Psi_Y:D^b_c(Y_{\bar\eta},E)
                                  \longrightarrow D^b_c(Y_s,E).
\]
The geometric generic fiber is part of this notation.  In particular,
$\Psi_Y$ takes no inertia invariants.  When an object has descent data,
nearby cycles may retain the resulting inertia action, but the underlying
complex is the one used here.

It is useful to name the natural exchange maps.  For a morphism
$f:Y\to Z$ of models, and complexes $A,B$ on geometric generic fibers, they
are
\begin{align}
 \operatorname{Ex}^*_f &: f_s^*\Psi_Z A
                  \longrightarrow\Psi_Y(f_{\bar\eta}^*A),
                          \label{eq:found-pull-exchange}\\
 \operatorname{Ex}_{f,*} &: \Psi_Z(f_{\bar\eta,*}B)
                  \longrightarrow f_{s,*}\Psi_Y B,
                          \label{eq:found-push-exchange}\\
 \mu_{A,B} &: \Psi_Y A\otimes\Psi_Y B
                  \longrightarrow\Psi_Y(A\otimes B).
                          \label{eq:found-tensor-exchange}
\end{align}
We use the push exchange only where it is defined by the proper
nearby-cycle comparison, and therefore an isomorphism.  Pull exchange and
tensor exchange are not asserted to be isomorphisms in general.  For two
models $Y,Z$, their composite gives the exterior-product map
\begin{equation}\label{eq:found-exterior-exchange}
 \operatorname{Ex}^{\boxtimes}_{Y,Z}:
  \Psi_Y A\boxtimes\Psi_Z B
       \longrightarrow\Psi_{Y\times_S Z}(A\boxtimes B).
\end{equation}
The products in this formula are on the appropriate fibers of the product
over $S$.

\begin{proposition}[Geometric nearby-cycle formalism]
\label{prop:nearby-foundations}
For the traits just specified, geometric nearby cycles preserve bounded
constructibility and are exact for the perverse $t$-structures on the
geometric fibers.  They are unchanged by finite extension of the trait
inside $\overline K$.  The maps
\eqref{eq:found-pull-exchange}, \eqref{eq:found-push-exchange}, and
\eqref{eq:found-exterior-exchange} are isomorphisms, respectively, for
smooth $f$, for proper $f$, and for exterior products of bounded
constructible complexes.  These assertions extend to locally
constructible complexes on the smooth stacks used here by smooth descent;
proper morphisms in this extension are representable.

Let $L_{\bar\eta}$ be a lisse finite-rank $E$-sheaf on a model's geometric
generic fiber.  Suppose it has, after a finite extension of the coefficient
field, a lisse lattice whose finite reductions extend to lisse sheaves
after finite extensions of $S$.  Suppose these extensions are compatible
on further common extensions and have reductions of a lisse special-fiber
sheaf $L_s$.  The finite extension of $S$ may depend on the reduction.
Then, for any map $g:Y\to Z$ to that model and any
$F\in D^b_c(Y_{\bar\eta},E)$, the natural map is an isomorphism
\begin{equation}\label{eq:found-lisse-tensor}
  \Psi_Y F\otimes g_s^*L_s
       \xrightarrow{\ \sim\ }
             \Psi_Y(F\otimes g_{\bar\eta}^*L_{\bar\eta}).
\end{equation}
The statement also holds locally on the stacks in question, and is
compatible with tensor products and morphisms of the lisse systems.
\end{proposition}

\begin{proof}
We recall the coefficient issue because it prevents an unwarranted descent
hypothesis in the applications.  At torsion coefficient level, a
constructible object on $Y_{\bar\eta}$ descends to a finite extension of
$K$; the same holds for any specified finite diagram of objects and maps.
Geometric nearby cycles of those descents agree after further finite
extensions.  Hence the torsion functor is defined on the filtered union of
these categories, independently of all such choices.  Applying the adic
formalism and extension of coefficients gives the geometric functor above.
This construction is the specialization functor of
\cite[\S\S4.1.1--4.1.5]{GR}.  Different reductions of an adic object can
require different extensions of $K$; the construction does not replace this
tower by a single finite extension.

The constructibility and exchange assertions are recorded in
\cite[\S4.1.6]{GR}.  For the general traits used here, they follow from
\cite[Theorem~4.1 and Corollary~4.2]{HansenScholze}, in the rational-adic
coefficient setting~(C) of that paper.  Indeed, the integral closure $V$ of
$R$ in $\overline K$ is a rank-one absolutely integrally closed valuation
ring.  On a separated finitely presented $V$-scheme, restriction to the
generic fiber identifies universally locally acyclic complexes with
arbitrary constructible generic-fiber complexes; its inverse is $Rj_*$.
Nearby cycles are $i^*Rj_*$.  This applies to $E$ itself, and thus directly
to objects without common finite descent.  Perverse exactness follows
from this equivalence and the relative perverse $t$-structure of
\cite[Theorem~6.7]{HansenScholze}: generic restriction and special
restriction are exact for their fiberwise perverse structures.  Finite
extensions of the original trait give the same $V$.  On a smooth stack,
smooth pullback of relative dimension $d$
commutes with $\Psi$ and with $[d]$.  The scheme assertions therefore
descend, and give both local constructibility and perverse exactness on
the stack.  The proper assertion is checked after smooth base change to
schemes, or algebraic spaces, for a representable proper map.

For the last assertion choose the lattice and reduce modulo a power of its
uniformizer.  After the permitted trait extension its extension is lisse
on $Z$.  The lisse tensor formula for nearby cycles, after pullback to
$Y$, gives \eqref{eq:found-lisse-tensor} at this finite coefficient level.
Its construction is natural in reductions, maps, and further trait
extensions.  Passing to the adic system and then inverting the coefficient
uniformizer proves the assertion.  This proof permits arbitrary $g$:
it uses lisseness of the extended factor, not a base-change theorem for
$\Psi$ along $g$.
\end{proof}

All exchange maps above obey the following compatibilities.  Pull and
proper-push exchange compose for composite morphisms; tensor exchange is
associative and unital; and these maps commute with proper base change and
the projection formula.  They are the natural transformations obtained
from restriction and pushforward in the definition of nearby cycles, so
these identities hold before choosing any isomorphisms.  This observation
will identify the maps obtained by different orders of Hecke convolution.

\subsection{Relative Satake normalization}

For a dominant coweight $\lambda$, put
$d_\lambda=\langle2\rho,\lambda\rangle$.  The spherical Satake complex
$\IC_\lambda$ is normalized on the open Schubert orbit by
\begin{equation}\label{eq:found-satake-normalization}
               \IC_\lambda|_{\Gr^\lambda}
                         =E[d_\lambda](d_\lambda/2).
\end{equation}
Since $G$ is simply connected, its coweight lattice is the coroot lattice;
$\langle2\rho,\alpha^\vee\rangle=2$ for every simple coroot.  Thus every
$d_\lambda$ is even.  We choose the usual Satake equivalence with
$\Rep(\Gd)$, including its cohomological fiber functor and fusion
commutativity constraint; see \cite[\S\S4--6 and Theorem~14.1]{MV}.
Tate twists are geometric coefficient lines.  They require no Weil
structure, and compatible trivializations may be fixed in comparisons
with Betti coefficients.

Let $U$ be the smooth scheme locus of the curve on which modifications are
allowed, disjoint from all marked or stacky points.  For a finite set $I$,
let $\cH_I$ be the Beilinson--Drinfeld Hecke stack with input bundle map
$p_I$ and output-and-leg map
\[
                       o_I:\cH_I\longrightarrow\cB\times U^I.
\]
For representations $\boldsymbol V=(V_i)_{i\in I}$, let
$\operatorname{Sat}_I(\boldsymbol V)$ denote its relative spherical
kernel.  In input frames at pairwise distinct legs it is the exterior
product of the fixed-point Satake complexes, with normalization
\eqref{eq:found-satake-normalization}.  At collisions it is obtained by
the proper convolution map from successive modifications.  Define
\begin{equation}\label{eq:found-hecke-definition}
 \Hecke_{I,\boldsymbol V}(F)
   =o_{I,*}\bigl(F\,\widetilde\boxtimes\,
                                  \operatorname{Sat}_I(\boldsymbol V)\bigr).
\end{equation}
The twisted product is characterized in input frames by the ordinary
exterior product with $F$.  The equivariance of the Satake factor provides
its descent.  These formulas are computed on finite-dimensional bounds
containing the support of the kernel; the output map on each such bound
is representable and proper.  No shift $[|I|]$ occurs in
\eqref{eq:found-hecke-definition}.  In particular
\begin{equation}\label{eq:found-hecke-unit}
 \Hecke_{I,(\mathbf1)}(F)=F\boxtimes E_{U^I},\qquad \Hecke_{\varnothing}(F)=F.
\end{equation}
When discussing perversity over a smooth leg space, we will explicitly
insert $[|I|]$ and subsequently remove it.

For a surjection of finite sets $a:I\twoheadrightarrow J$, write
$\delta_a:U^J\to U^I$ for its collision diagonal, and put
$W_j=\bigotimes_{i\in a^{-1}(j)}V_i$.  The relative fusion convention is
\begin{equation}\label{eq:found-fusion}
  (1\times\delta_a)^*\Hecke_{I,\boldsymbol V}(F)
                         \simeq \Hecke_{J,\boldsymbol W}(F).
\end{equation}
There is no codimension shift in this ordinary-pullback formula.  The
usual shifts appear only if both sides have first been made perverse over
their respective leg spaces.  An eigenstructure throughout this paper
means natural isomorphisms for all $I$ and all representations, compatible
with these fusion maps, the unit, successive modifications, and
permutations.  A system of separate fixed-point eigenisomorphisms is not
sufficient.

\subsection{The use of Betti realization}

For a finite-type complex scheme, geometric adic constructible complexes
admit the usual realization as algebraically constructible complexes of
ordinary $E$-vector spaces on the analytic space.  It is obtained from
finite coefficient comparison and passage to adic coefficients; on a
lisse sheaf it restricts its continuous monodromy representation to
topological loops.  On bounded constructible objects the realization is
conservative, preserves the perverse $t$-structure, and commutes with the
pullbacks, tensor products, and finite-type pushforwards used below.  The
local finite-triangulation form of comparison ensures that these
operations commute with passage from a lattice to $E$-coefficients.
The characteristic-zero comparison convention is also used in
\cite[\S2.2.1]{GR}.

For algebraically constructible complexes this comparison identifies
algebraic singular support with the corresponding complex conic singular
support in the analytic cotangent bundle.  One can check the assertion
using the vanishing-cycle characterization: comparison identifies the
tests by algebraic functions, and such tests generate singular support.
We apply these statements on finite-type smooth charts and their
correspondence diagrams, and descend them to the stacks above.  In
particular Betti realization may verify that an already defined adic
truncation comparison is an isomorphism.  It does not create an adic
object from an arbitrary Betti object.  The infinite-rank universal
monodromic kernels appearing later are used only on the Betti side; all
objects subsequently specialized are geometric adic lcc complexes.
